
\documentclass{article} % For LaTeX2e
\usepackage{iclr2027_conference,times}

% Optional math commands from https://github.com/goodfeli/dlbook_notation.
\input{math_commands.tex}

\usepackage{hyperref}
\usepackage{url}
 

\title{HomeACE: A Comprehensive Evaluation Benchmark for Home Anomaly Video Understanding}

% Authors must not appear in the submitted version. They should be hidden
% as long as the \iclrfinalcopy macro remains commented out below.
% Non-anonymous submissions will be rejected without review.

\author{Antiquus S.~Hippocampus, Natalia Cerebro \& Amelie P. Amygdale \thanks{ Use footnote for providing further information
about author (webpage, alternative address)---\emph{not} for acknowledging
funding agencies.  Funding acknowledgements go at the end of the paper.} \\
Department of Computer Science\\
Cranberry-Lemon University\\
Pittsburgh, PA 15213, USA \\
\texttt{\{hippo,brain,jen\}@cs.cranberry-lemon.edu} \\
\And
Ji Q. Ren \& Yevgeny LeNet \\
Department of Computational Neuroscience \\
University of the Witwatersrand \\
Joburg, South Africa \\
\texttt{\{robot,net\}@wits.ac.za} \\
\AND
Coauthor \\
Affiliation \\
Address \\
\texttt{email}
}

% The \author macro works with any number of authors. There are two commands
% used to separate the names and addresses of multiple authors: \And and \AND.
%
% Using \And between authors leaves it to \LaTeX{} to determine where to break
% the lines. Using \AND forces a linebreak at that point. So, if \LaTeX{}
% puts 3 of 4 authors names on the first line, and the last on the second
% line, try using \AND instead of \And before the third author name.

\newcommand{\fix}{\marginpar{FIX}}
\newcommand{\new}{\marginpar{NEW}}

%\iclrfinalcopy % Uncomment for camera-ready version, but NOT for submission.
\begin{document}


\maketitle

\begin{abstract}
Understanding potential risks and anomalous behaviors involving household members in real-world home environments requires models to detect subtle anomalies amid complex and changing surroundings, integrate long-term temporal context, and maintain reliable perception and reasoning. This capability is a critical measure of whether large multimodal models can support video analysis in domestic settings and be deployed in home-monitoring cameras. However, existing evaluations of large multimodal models mainly focus on general video understanding, visual question answering (VQA), temporal localization, and related tasks, with limited attention to risks and anomalies in home environments and no systematic assessment of the comprehensive reasoning abilities required for this setting. To address this gap, we introduce \textbf{HomeACE}, a large-scale \textbf{Home} \textbf{A}nomaly \textbf{C}omprehensive \textbf{E}valuation dataset for home anomaly understanding. The HomeACE dataset comprises more than 16K videos of domestic scenes, including nearly 3K real-world videos and 13K synthetic videos, covering anomalous, risky, and normal situations. Based on this dataset, we develop HomeACE-Bench, a dedicated benchmark that evaluates models across four complementary dimensions: Perception, Cognition, Temporal reasoning, and Planning. Extensive quantitative and qualitative experiments show that state-of-the-art multimodal models exhibit substantial deficiencies in anomaly perception, interference robustness under complex household contexts, and temporal reasoning, with significant room for improvement in generalization to real-world homes. The resulting HomeACE suite provides a standardized evaluation platform for advancing home anomaly understanding models and fostering the development of robust and practical intelligent anomaly-aware systems for real-world domestic applications.
\end{abstract}

\section{Introduction}

Recent advances in multimodal large language models (MLLMs) have greatly accelerated the deployment of video understanding technology, providing important technical support for downstream applications such as smart-home monitoring, domestic safety, and intelligent perception in household environments. Understanding anomalous and risky behaviors in real homes requires models to accurately perceive risky actions and potential safety hazards in complex and changing environments, infer and describe their causes and complete development processes, precisely localize the temporal intervals in which they occur, and produce reasonable risk-mitigation plans to prevent similar incidents from recurring. These capabilities are core criteria for determining whether MLLMs can adapt to video analysis in home environments and be deployed in consumer home-monitoring devices. They also represent a key bottleneck for the large-scale deployment of intelligent home-safety systems.

However, existing evaluation frameworks for multimodal models mainly focus on general video understanding, visual question answering (VQA), temporal localization, and other foundational tasks. Benchmarks built around generic scenes do not adequately capture the distinctive properties of household environments and cannot comprehensively evaluate the four capabilities required for home-safety understanding: risk perception, process reasoning, temporal localization, and strategy planning. Meanwhile, existing public datasets generally suffer from limited scene coverage. They cannot encompass the rich diversity of household environments or reproduce the stochastic and heterogeneous nature of anomalous and risky behaviors in homes, severely limiting the effective evaluation of models' understanding of real-world domestic scenarios.

To address this gap, we construct \textbf{HomeACE}, a large-scale and comprehensive dataset for home anomaly and risk understanding. HomeACE contains more than 16K videos of household scenes, including nearly 3K real-world videos and 13K synthetic videos. It covers routine household activities, a broad range of anomalous behaviors, and potential safety-risk scenarios, substantially exceeding existing benchmarks in both scale and scene diversity.

Building on this high-quality dataset, we further develop HomeACE-Bench, a dedicated evaluation benchmark centered on the core requirements of home anomaly understanding. It evaluates model capabilities along four complementary and comprehensive dimensions: Perception, Cognition, Temporal reasoning, and Planning. The Perception dimension measures a model's ability to identify explicit and implicit anomalies and potential risks in household scenes. Cognition focuses on understanding the causes of risks and inferring their unfolding processes. Temporal reasoning evaluates the precise localization of when a risk begins, persists, and ends. Planning measures the ability to formulate reasonable risk-mitigation strategies and prevent the recurrence of similar risks. By covering the core functional requirements of home-safety monitoring, this evaluation framework moves beyond the limitations of traditional benchmarks based primarily on isolated detection and classification tasks, providing a more realistic assessment of models' performance in practical deployments.

Extensive quantitative and qualitative experiments conducted with HomeACE-Bench reveal substantial deficiencies in current state-of-the-art MLLMs for household scenarios. These models generally exhibit limited accuracy in identifying anomalous risks, weak reasoning about the causes and development processes of risks, imprecise temporal localization, and incomplete risk-mitigation plans. Their generalization to unconstrained real-world home environments also remains poor, making them inadequate for the practical requirements of home-safety monitoring.

In summary, HomeACE and HomeACE-Bench provide a standardized and comprehensive evaluation platform for home anomaly understanding. They can drive the iterative improvement of models in this domain and facilitate the development of robust, practical, and deployable intelligent anomaly-perception systems for domestic environments. Our main contributions are summarized as follows:

\begin{enumerate}
    \item \textbf{Large-scale household-scene dataset.} We construct HomeACE, a large-scale dataset for home anomaly understanding that combines real-world and synthetic scenes. It contains more than 16K household videos covering routine, anomalous, and risky situations, addressing the limitations of existing datasets in scale, real-world coverage, and scene diversity.
    \item \textbf{Multi-dimensional systematic evaluation benchmark.} We propose HomeACE-Bench, a dedicated evaluation framework that establishes a comprehensive paradigm across four dimensions: anomaly perception, risk cognition, temporal localization, and solution planning. It overcomes the narrow focus of conventional evaluation tasks and enables fine-grained, systematic diagnosis of models' overall capabilities in household scenarios.
    \item \textbf{Systematic analysis of model limitations.} Through extensive benchmark experiments, we conduct a detailed analysis of the key deficiencies of state-of-the-art MLLMs across the four core capabilities of home anomaly understanding. We perform targeted studies from multiple perspectives, including data-type differences, anomalous semantic understanding, and the temporal distribution of anomalies, and investigate how factors such as video duration and temporal position affect model performance. These analyses identify the core bottlenecks in home risk perception, reasoning, localization, and planning, providing concrete guidance for improving the next generation of intelligent domestic-perception models.
\end{enumerate}


\section{Related Work}
\subsection{Multimodal Large Language Models for Video Understanding}

Multimodal large language models (MLLMs) have rapidly expanded video understanding from frame-level recognition to open-ended video reasoning and dialogue. Early video MLLMs typically adapt an image-language architecture by sampling video frames and aligning visual features with a language model. Video-ChatGPT, for example, combines a video-adapted visual encoder with an LLM and introduces video instruction tuning for detailed video conversations \citep{Maaz2024VideoChatGPT}. Subsequent models improve the modeling of spatiotemporal structure and multimodal cues. VideoLLaMA 2 incorporates a spatial-temporal connector and an audio branch to support video, audio, and audio-visual question answering \citep{Cheng2024VideoLLaMA2}, while LLaVA-OneVision demonstrates transfer across single-image, multi-image, and video scenarios within a unified open MLLM \citep{Li2024LLaVAOneVision}. More recent efforts have turned to long-context video reasoning, emphasizing continuous tracking, information integration, and memory retention over extended temporal spans. VideoOdyssey, for instance, evaluates visual and audio-visual understanding under ultra-long videos and multiple levels of continuous evidence requirements \citep{He2026VideoOdyssey}. Despite these advances, existing MLLMs and their evaluations are largely developed around generic video content and conventional understanding tasks. They do not specifically test whether a model can identify subtle risks in a home, explain the causal development of an incident, localize its complete temporal interval, and formulate an actionable mitigation strategy. HomeACE complements this line of research by evaluating these capabilities in the structured yet unconstrained context of real-world domestic environments.

\subsection{Home Video Benchmarks}

Existing home-video datasets cover several important but relatively narrow aspects of household understanding. Toyota Smarthome was recorded in an apartment with multiple Kinect cameras and contains daily living activities performed by older adults. Its trimmed version supports classification of 31 activities with 16,115 short RGB-D samples, whereas Toyota Smarthome Untrimmed retains long recordings and densely annotates 51 activities for temporal activity detection \citep{Das2019ToyotaSmarthome,Dai2022ToyotaSmarthomeUntrimmed}. YouHome introduces a home-oriented system and dataset for modeling everyday activities and contextual information in domestic environments \citep{Pan2022YouHome}. These datasets provide realistic settings and valuable benchmarks for activity recognition and detection, but they primarily focus on predefined daily activities rather than open-ended anomalous behaviors, potential safety risks, or risk-mitigation reasoning.

Another group of datasets targets specific safety events, especially human falls. The UR Fall Detection Dataset contains 70 sequences, including 30 falls and 40 activities of daily living, recorded with Kinect cameras and accelerometric sensors \citep{Kwolek2014URFall}. The Multiple Cameras Fall Dataset focuses on multi-camera fall detection and provides a testbed for handling viewpoint changes and occlusions in surveillance settings \citep{Auvinet2008MultipleCameraFall}. While these datasets are important for developing fall-detection systems, their task definitions are centered on detecting a particular event and do not cover the broader spectrum of household anomalies or the reasoning and planning capabilities required for comprehensive home-safety understanding.

More recently, SmartHome-Bench was proposed as a benchmark for video anomaly detection in smart-home scenarios using MLLMs \citep{Zhao2025SmartHomeBench}. It contains 1,203 smart-home-camera videos with anomaly tags, descriptions, and reasoning annotations, and organizes anomalies into several application-oriented categories such as wildlife, senior care, and baby monitoring. SmartHome-Bench moves beyond general-purpose video anomaly benchmarks by considering the characteristics of smart-home monitoring. Nevertheless, it remains primarily focused on anomaly detection and does not jointly evaluate anomaly perception, causal and process cognition, precise temporal reasoning, and risk-mitigation planning. Overall, existing home-video resources are fragmented across daily-activity recognition, fall detection, and anomaly detection, with limited coverage of diverse household contexts and limited support for systematic evaluation of comprehensive home anomaly understanding. HomeACE addresses these limitations by combining real-world and synthetic household videos with a unified benchmark spanning Perception, Cognition, Temporal reasoning, and Planning.


\section{HomeACE}


\section{Final instructions}
Do not change any aspects of the formatting parameters in the style files.
In particular, do not modify the width or length of the rectangle the text
should fit into, and do not change font sizes (except perhaps in the
\textsc{References} section; see below). Please note that pages should be
numbered.

\section{Preparing PostScript or PDF files}

Please prepare PostScript or PDF files with paper size ``US Letter'', and
not, for example, ``A4''. The -t
letter option on dvips will produce US Letter files.

Consider directly generating PDF files using \verb+pdflatex+
(especially if you are a MiKTeX user).
PDF figures must be substituted for EPS figures, however.

Otherwise, please generate your PostScript and PDF files with the following commands:
\begin{verbatim}
dvips mypaper.dvi -t letter -Ppdf -G0 -o mypaper.ps
ps2pdf mypaper.ps mypaper.pdf
\end{verbatim}

\subsection{Margins in LaTeX}

Most of the margin problems come from figures positioned by hand using
\verb+\special+ or other commands. We suggest using the command
\verb+\includegraphics+
from the graphicx package. Always specify the figure width as a multiple of
the line width as in the example below using .eps graphics
\begin{verbatim}
   \usepackage[dvips]{graphicx} ...
   \includegraphics[width=0.8\linewidth]{myfile.eps}
\end{verbatim}
or % Apr 2009 addition
\begin{verbatim}
   \usepackage[pdftex]{graphicx} ...
   \includegraphics[width=0.8\linewidth]{myfile.pdf}
\end{verbatim}
for .pdf graphics.
See section~4.4 in the graphics bundle documentation (\url{http://www.ctan.org/tex-archive/macros/latex/required/graphics/grfguide.ps})

A number of width problems arise when LaTeX cannot properly hyphenate a
line. Please give LaTeX hyphenation hints using the \verb+\-+ command.

\subsection*{AI use statement}

(This section is \textbf{required} and does not count toward the page limit.)

In this work, we used generative AI tools for [tasks with required disclosure].
We have not used generative AI tools for [other tasks with required disclosure],
and [the rest of the required disclosure tasks] are not applicable to this work.
Additionally, we used generative AI tools for [tasks with recommended
disclosure]. We have reviewed all AI-assisted work. [Elaborate. For example, “we
checked LLM-generated research ideas for potential plagiarism through a manual
literature survey”, “LLM-generated code was verified and tested for correctness
by 2 authors”, etc.]. We take responsibility for the final content of this work,
including text, claims or artifacts produced with the aid of generative AI.

See the ICLR 2027 AI Policy for Authors for more details. This statement should
not be more than 1 page.

\subsection*{Ethics statement}

(This section is \textbf{recommended} and does not count toward the page limit.)

If authors feel that their paper submission raises questions regarding the Code
of Ethics, they are encouraged to include a paragraph of Ethics Statement (at
the end of the main text before references) to address potential concerns where
appropriate. Topics include, but are not limited to, studies that involve human
subjects, practices to data set releases, potentially harmful insights,
methodologies and applications, potential conflicts of interest and sponsorship,
discrimination/bias/fairness concerns, privacy and security issues, legal
compliance, and research integrity issues (e.g., IRB, documentation, research
ethics). This statement should not be more than 1 page.

\subsection*{Reproducibility statement}

(This section is \textbf{recommended} and does not count toward the page limit.)

It is important that the work published in ICLR is reproducible. Authors are
strongly encouraged to include a paragraph-long Reproducibility Statement at the
end of the main text (before references) to discuss the efforts that have been
made to ensure reproducibility. This paragraph should not itself describe
details needed for reproducing the results, but rather reference the parts of
the main paper, appendix, and supplemental materials that will help with
reproducibility. For example, for novel models or algorithms, a link to an
anonymous downloadable source code can be submitted as supplementary materials;
for theoretical results, clear explanations of any assumptions and a complete
proof of the claims can be included in the appendix; for any datasets used in
the experiments, a complete description of the data processing steps can be
provided in the supplementary materials. Each of the above are examples of
things that can be referenced in the reproducibility statement.

\subsubsection*{Author Contributions}
If you'd like to, you may include  a section for author contributions as is done
in many journals. This is optional and at the discretion of the authors.

\subsubsection*{Acknowledgments}
Use unnumbered third level headings for the acknowledgments. All
acknowledgments, including those to funding agencies, go at the end of the paper.


\bibliography{iclr2027_conference}
\bibliographystyle{iclr2027_conference}

\appendix
\section{Appendix}
You may include other additional sections here.


\end{document}
